\chapter{De voortplantingsstelsels}\label{ch:g8:reproductive-systems}

\Cref{ch:g8:sexual-reproduction} had \omterm{def:g8:sexual-reproduction:gametes}{gameten} nodig;
\cref{ch:g8:puberty} zette hun fabrieken aan. Dit hoofdstuk brengt de
fabrieken zelf in kaart: het mannelijke en het vrouwelijke
voortplantingsstelsel --- de organen, de \omterm{def:g8:sexual-reproduction:gametes}{gameten} die elk maakt, en het
opmerkelijke maandelijkse ritme van het vrouwelijke stelsel. Gewone
anatomie, gewoon benoemd: dit zijn organen zoals \omterm{def:g4:heart-and-blood:heart}{hart} en \omterm{def:g7:how-animals-breathe:four}{longen}, en dit is
hun hoofdstuk.

\section{Het mannelijke stelsel}

\begin{definition}[Het mannelijke voortplantingsstelsel]\label{def:g8:reproductive-systems:male}
De organen van het mannelijke stelsel:
\begin{itemize}
  \item de twee \emph{zaadballen}\index{zaadbal}, buiten de warmte van het
        lichaam in de \emph{balzak} gehouden --- de gametenfabrieken:
        vanaf de \omterm{def:g8:puberty:puberty}{puberteit} maken ze onophoudelijk
        \emph{zaadcellen}\index{zaadcel}, bij honderden miljoenen;
  \item buizen die het \omterm{def:g2:from-seed-to-plant:seed}{zaad} opslaan en vervoeren, onderweg vergezeld door
        klieren waarvan het vocht de zaadcellen voedt en meevoert ---
        samen vormen ze het \emph{sperma};
  \item de \emph{penis}\index{penis}, waardoor het sperma het lichaam
        verlaat --- hetzelfde orgaan dat op zijn andere dienst \omterm{def:g7:removing-wastes:kidneys}{urine}
        afvoert (\cref{def:g7:removing-wastes:kidneys}); een stelsel van
        \omterm{def:g7:heart-and-circulation:rooms}{kleppen} houdt de twee diensten strikt gescheiden.
\end{itemize}
\end{definition}

\begin{example}[De zaadcel, van dichtbij]\label{ex:g8:reproductive-systems:sperm}
Onder de \omterm{def:g5:first-look-at-cells:microscope}{microscoop}
(\cref{def:g5:first-look-at-cells:microscope}) is een \omterm{def:g8:reproductive-systems:male}{zaadcel} een \omterm{def:g5:first-look-at-cells:cell}{cel} die
tot op het \omterm{def:g3:bones-and-muscles:skeleton}{bot} voor de reis is uitgekleed: een kop die bijna helemaal uit
\omterm{def:g6:cell-unit-of-life:parts}{celkern} bestaat --- de ingepakte bijdrage van de vader
(\cref{rem:g5:first-look-at-cells:nucleus}) --- en een zwiepende \omterm{def:g1:animals-around-us:parts}{staart}.
De kleinste \omterm{def:g5:first-look-at-cells:cell}{cel} van het lichaam, gebouwd om te zwemmen, voor een paar
dagen van proviand voorzien. De productie loopt van de \omterm{def:g8:puberty:puberty}{puberteit} tot in de
\omterm{def:g3:stages-of-life:stages}{ouderdom}: de mannelijke fabriek draait onafgebroken.
\end{example}

\section{Het vrouwelijke stelsel}

\begin{definition}[Het vrouwelijke voortplantingsstelsel]\label{def:g8:reproductive-systems:female}
De organen van het vrouwelijke stelsel, alle in de onderbuik:
\begin{itemize}
  \item de twee \emph{eierstokken}\index{eierstok} --- de voorraden en
        fabrieken van \omterm{def:g8:sexual-reproduction:gametes}{gameten}: anders dan de \omterm{def:g8:reproductive-systems:male}{zaadballen} dragen ze hun
        voorraad toekomstige \emph{\omterm{def:g4:animal-reproduction:sexual}{eicellen}} al vanaf de \omterm{def:g2:born-grow-die:cycle}{geboorte}, en
        vanaf de \omterm{def:g8:puberty:puberty}{puberteit} laten ze er één tegelijk los;
  \item de twee \emph{eileiders}\index{eileider}, met franjes afgezette
        buizen die de losgelaten \omterm{def:g4:animal-reproduction:sexual}{eicel} opvangen --- de gebruikelijke
        ontmoetingsplaats van stap 2 uit
        \cref{prop:g8:sexual-reproduction:core};
  \item de \emph{\omterm{prop:g5:human-reproduction:plan}{baarmoeder}} (de \omterm{prop:g5:human-reproduction:plan}{baarmoeder} uit
        \cref{prop:g5:human-reproduction:plan}), waarvan de dikke
        spierwand en het vernieuwbare slijmvlies het onderdak van de
        \omterm{prop:g8:sexual-reproduction:core}{zygote} vormen;
  \item de \emph{vagina}, die de \omterm{prop:g5:human-reproduction:plan}{baarmoeder} met de buitenwereld verbindt
        --- waar het sperma ontvangen wordt, en het geboortekanaal uit
        \cref{ex:g5:human-reproduction:birth}.
\end{itemize}
\end{definition}

\begin{example}[De eicel, van dichtbij]\label{ex:g8:reproductive-systems:egg}
De \omterm{def:g4:animal-reproduction:sexual}{eicel} is in alles behalve haar rol het tegendeel van de \omterm{def:g8:reproductive-systems:male}{zaadcel}: de
\emph{grootste} \omterm{def:g5:first-look-at-cells:cell}{cel} van het lichaam --- een tiende millimeter, net aan de
grens van het zichtbare --- rond, onbeweeglijk, en voorzien van proviand
voor de eerste dagen van de reis (het dooierprincipe uit
\cref{rem:g2:animals-grow-up:egg}, op celschaal). De ene fabriek draait
onafgebroken miljoenen; de andere laat één van proviand voorziene \omterm{def:g5:first-look-at-cells:cell}{cel} per
maand los: het aantallen-tegenover-zorg uit
\cref{prop:g4:animal-reproduction:strategies}, in de \omterm{def:g8:sexual-reproduction:gametes}{gameten} zelf
geschreven.
\end{example}

\section{De cyclus}

\begin{proposition}[De menstruatiecyclus]\label{prop:g8:reproductive-systems:cycle}
Van de \omterm{def:g8:puberty:puberty}{puberteit} tot ongeveer het vijftigste jaar draait het vrouwelijke
stelsel een zich herhalend programma van ruwweg \emph{28 dagen} --- de
\emph{menstruatiecyclus}\index{menstruatiecyclus}, bevolen door de
hormoonketen uit \cref{def:g8:puberty:hormones}:
\begin{enumerate}
  \item het slijmvlies van de \omterm{prop:g5:human-reproduction:plan}{baarmoeder} bouwt zichzelf op, dik en rijk
        aan \omterm{prop:g4:journey-of-food:absorb}{bloed} --- het onderdak wordt klaargemaakt;
  \item halverwege de cyclus, rond dag 14, laat een \omterm{def:g8:reproductive-systems:female}{eierstok} één \omterm{def:g4:animal-reproduction:sexual}{eicel} los
        --- de \emph{eisprong}\index{eisprong}; de \omterm{def:g8:reproductive-systems:female}{eileider} vangt haar op,
        en ongeveer een dag lang kan ze bevrucht worden;
  \item geen \omterm{def:g5:flower-to-fruit:steps}{bevruchting}: het slijmvlies is niet nodig; het breekt af en
        verlaat het lichaam in een paar dagen via de vagina --- de
        \emph{menstruatie}\index{menstruatie} --- en het programma begint
        opnieuw;
  \item wel \omterm{def:g5:flower-to-fruit:steps}{bevruchting}: de \omterm{prop:g8:sexual-reproduction:core}{zygote} nestelt zich in het klaargemaakte
        slijmvlies
        (\cref{def:g5:human-reproduction:pregnancy}), de cyclus wordt
        opgeschort, en de \omterm{def:g5:human-reproduction:pregnancy}{zwangerschap} begint.
\end{enumerate}
\end{proposition}
\admitted

\begin{omfigure}
\begin{tikzpicture}[scale=0.94]
  % cycle wheel
  \draw[very thick, omDef!60] (0,0) circle (2.3);
  % day marks
  \foreach \a/\d in {90/1, 0/8, -90/14, 180/22} {
    \fill[omDef] (\a:2.3) circle (0.07);
  }
  \node[font=\footnotesize, above] at (90:2.45) {dag 1};
  \node[font=\footnotesize, right] at (0:2.45) {dag 8};
  \node[font=\footnotesize, below] at (-90:2.45) {dag 14};
  \node[font=\footnotesize, left] at (180:2.45) {dag 22};
  % phases
  \draw[very thick, omThm] (90:2.3) arc[start angle=90, end angle=35,
    radius=2.3];
  \node[font=\footnotesize, omThm, align=center] at (58:3.5)
    {de menstruatie:\\ het slijmvlies gaat eraf};
  \node[font=\footnotesize, omMeth, align=center] at (-35:4.0)
    {slijmvlies opgebouwd,\\ dik en klaar};
  \fill[omProp] (-90:2.3) circle (0.12);
  \node[font=\footnotesize, omProp, align=center] at (-90:3.55)
    {eisprong:\\ één eicel losgelaten};
  % center
  \node[font=\small, align=center] at (0,0.25)
    {de cyclus:\\ ongeveer 28 dagen};
  \node[font=\footnotesize, align=center] at (0,-0.75)
    {(bevruchting schort haar op:\\ zwangerschap)};
  % arrow direction
  \draw[->, thick, omDef] (150:2.3) arc[start angle=150,
    end angle=120, radius=2.3];
\end{tikzpicture}

{\small De \omterm{prop:g8:reproductive-systems:cycle}{menstruatiecyclus} als een wiel: de \omterm{prop:g8:reproductive-systems:cycle}{menstruatie} opent hem, het
slijmvlies bouwt zich op, de \omterm{prop:g8:reproductive-systems:cycle}{eisprong} bekroont hem halverwege --- en de
\omterm{def:g5:flower-to-fruit:steps}{bevruchting} is het ene voorval dat het wiel stilzet.}
\end{omfigure}

\begin{remark}[Leven met de cyclus]\label{rem:g8:reproductive-systems:living}
Praktische zaken, gewoon gezegd: de eerste jaren zijn de cycli
onregelmatig --- de machinerie stelt zich af; de lengte verschilt van
persoon tot persoon en van maand tot maand; \omterm{prop:g8:reproductive-systems:cycle}{menstruaties} kunnen krampen
geven (de \omterm{prop:g5:human-reproduction:plan}{baarmoeder} is een \omterm{def:g3:bones-and-muscles:muscle}{spier},
\cref{def:g3:bones-and-muscles:muscle}, en die werkt terwijl het
slijmvlies afgestoten wordt) --- warmte, beweging en, als het nodig is,
het advies van de schoolverpleegkundige helpen allemaal; en maandverband,
tampons en cups zijn gewoon de praktische uitrusting van de \omterm{prop:g8:reproductive-systems:cycle}{menstruatie},
een normaal schap in elke drogisterij. Niets daarvan is ziekte: de cyclus
is het stelsel dat precies draait zoals het gebouwd is.
\end{remark}

\begin{method}[De twee stelsels naast elkaar lezen]\label{met:g8:reproductive-systems:sides}
Om de kaart helder te houden, vergelijk je op functie:
\begin{enumerate}
  \item gametenfabriek: \omterm{def:g8:reproductive-systems:male}{zaadballen} --- \omterm{def:g8:reproductive-systems:female}{eierstokken};
  \item \omterm{def:g8:sexual-reproduction:gametes}{gameet}: onafgebroken zwemmende miljoenen --- één van proviand
        voorziene \omterm{def:g4:animal-reproduction:sexual}{eicel} per maand;
  \item ontmoetingsplaats: de \omterm{def:g8:reproductive-systems:female}{eileider}, bereikt door de zwemtocht van de
        \omterm{def:g8:reproductive-systems:male}{zaadcel} vanuit de vagina door de \omterm{prop:g5:human-reproduction:plan}{baarmoeder};
  \item onderdak: geen op de mannelijke kaart --- de \omterm{prop:g5:human-reproduction:plan}{baarmoeder}, maandelijks
        klaargemaakt, op de vrouwelijke;
  \item tijdschema: onafgebroken vanaf de \omterm{def:g8:puberty:puberty}{puberteit} --- cyclisch, van de
        \omterm{def:g8:puberty:puberty}{puberteit} tot ongeveer het vijftigste jaar.
\end{enumerate}
\end{method}

\section{Oefeningen}

\begin{exercise}[$\star$]\label{exo:g8:reproductive-systems:1}
Breng het mannelijke stelsel in kaart: noem elk orgaan en zijn functie.
\end{exercise}

\begin{exercise}[$\star$]\label{exo:g8:reproductive-systems:2}
Breng het vrouwelijke stelsel op dezelfde manier in kaart.
\end{exercise}

\begin{exercise}[$\star$]\label{exo:g8:reproductive-systems:3}
Zet \omterm{def:g8:reproductive-systems:male}{zaadcel} en \omterm{def:g4:animal-reproduction:sexual}{eicel} tegenover elkaar: grootte, bouw, aantallen,
proviand.
\end{exercise}

\begin{exercise}[$\star$]\label{exo:g8:reproductive-systems:4}
Wat gebeurt er bij de \omterm{prop:g8:reproductive-systems:cycle}{eisprong}, op welke dag ruwweg, en hoe lang is de
\omterm{def:g4:animal-reproduction:sexual}{eicel} te bevruchten?
\end{exercise}

\begin{exercise}[$\star$]\label{exo:g8:reproductive-systems:5}
Wat is de \omterm{prop:g8:reproductive-systems:cycle}{menstruatie}, in termen van het slijmvlies --- en wat zegt haar
komst over die cyclus?
\end{exercise}

\begin{exercise}[$\star$]\label{exo:g8:reproductive-systems:6}
Welk voorval schort de cyclus op, en wat volgt er in de \omterm{prop:g5:human-reproduction:plan}{baarmoeder}?
\end{exercise}

\begin{exercise}[$\star\star$]\label{exo:g8:reproductive-systems:7}
Waar gebeuren de drie stappen uit
\cref{prop:g8:sexual-reproduction:core} op de kaart van dit hoofdstuk ---
orgaan voor orgaan?
\end{exercise}

\begin{exercise}[$\star\star$]\label{exo:g8:reproductive-systems:8}
Lees de \omterm{def:g8:sexual-reproduction:gametes}{gameten} als de twee strategieën uit
\cref{prop:g4:animal-reproduction:strategies} in het klein. Welke kant
draagt welke logica?
\end{exercise}

\begin{exercise}[$\star\star$]\label{exo:g8:reproductive-systems:9}
Waarom zijn de eerste jaren cycli vaak onregelmatig, en waarom is een
wisselende lengte normaal? (Twee feiten uit
\cref{rem:g8:reproductive-systems:living}.)
\end{exercise}

\begin{exercise}[$\star\star$]\label{exo:g8:reproductive-systems:10}
Verklaar menstruatiekrampen met
\cref{def:g3:bones-and-muscles:muscle} --- en noem twee dingen die helpen.
\end{exercise}

\begin{exercise}[$\star\star$]\label{exo:g8:reproductive-systems:11}
De twee fabrieken draaien tegenovergestelde tijdschema's ---
onafgebroken tegenover cyclisch-met-voorraad-vanaf-de-geboorte. Verbind
elk met de bouw en de rol van haar \omterm{def:g8:sexual-reproduction:gametes}{gameet}.
\end{exercise}

\begin{exercise}[$\star\star\star$]\label{exo:g8:reproductive-systems:12}
``De cyclus is de \omterm{prop:g5:human-reproduction:plan}{baarmoeder} die elke maand de vraag stelt die alleen een
\omterm{def:g5:flower-to-fruit:steps}{bevruchting} kan beantwoorden.'' Vouw dat beeld open tot de vier genummerde
gebeurtenissen uit
\cref{prop:g8:reproductive-systems:cycle}.
\end{exercise}

\section{Probleem: De lesmodellen van de kliniek}

\begin{problem}[{Weekendprobleem --- de stelsels uitleggen met twee
modellen en een kalender}]\label{pb:g8:reproductive-systems:1}
Een gezondheidscentrum gebruikt op zijn open dag drie leermiddelen: een
model van het mannelijke stelsel, een model van het vrouwelijke stelsel,
en een grote kalender van 28 dagen. Jij bent de vrijwillige uitlegger.

\medskip\noindent\textbf{Deel I --- De twee modellen.}
\begin{enumerate}
  \item Een bezoeker volgt het mannelijke model van de \omterm{def:g8:reproductive-systems:male}{zaadballen} tot de
        uitgang. Vertel de route na, en noem wat er onderweg bijkomt en
        hoe het dubbeldienstorgaan geregeld is.
  \item Wijs op het vrouwelijke model aan waar de gametenvoorraad wacht,
        waar de \omterm{def:g5:flower-to-fruit:steps}{bevruchting} meestal gebeurt, en waar een \omterm{def:g5:human-reproduction:pregnancy}{zwangerschap}
        onderdak vindt.
  \item Een bezoeker vraagt waarom de gametenfabrieken van de modellen zo
        scherp verschillen --- miljoenen per dag tegenover één per maand.
        Antwoord met de slotzin van
        \cref{ex:g8:reproductive-systems:egg}.
  \item Iemand anders vraagt wat de stelsels vóór de \omterm{def:g8:puberty:puberty}{puberteit} deden.
        Antwoord in één regel met de machinerie van \cref{ch:g8:puberty}.
\end{enumerate}

\medskip\noindent\textbf{Deel II --- De kalender.}
\begin{enumerate}[resume]
  \item Loop de kalender langs: plaats de \omterm{prop:g8:reproductive-systems:cycle}{menstruatie}, de opbouw, de
        \omterm{prop:g8:reproductive-systems:cycle}{eisprong} en de vruchtbare dag, met hun globale data.
  \item Een bezoeker neemt aan dat dag 28 voor iedereen hetzelfde is.
        Verbeter die aanname met
        \cref{rem:g8:reproductive-systems:living} --- twee vormen van
        variatie.
  \item Laat op de kalender zien wat er verandert als er rond dag 14 een
        \omterm{def:g5:flower-to-fruit:steps}{bevruchting} plaatsvindt --- welke gebeurtenissen vervallen, welke
        beginnen
        (\cref{def:g5:human-reproduction:pregnancy} vertelt het verhaal
        verder).
  \item Een tiener vraagt zachtjes of krampen betekenen dat er iets mis
        is. Geef het spierantwoord en de twee hulpmiddelen --- en het
        adres voor alles wat verder gaat.
\end{enumerate}

\medskip\noindent\textbf{Deel III --- Het vak van de
uitlegger.}
\begin{enumerate}[resume]
  \item Bezoekers giechelen om de modellen. Jouw openingszin behandelt de
        stelsels zoals de eerste alinea van
        \cref{ch:g8:reproductive-systems} dat doet. Schrijf hem op.
  \item Vat de twee stelsels samen in de vergelijking van vijf regels uit
        \cref{met:g8:reproductive-systems:sides}, uit je \omterm{def:g1:my-body:parts}{hoofd}.
  \item Een bezoeker vraagt waar het begin van \omterm{def:g3:stages-of-life:stages}{baby}'s op de leermiddelen
        past. Rijg de drie leermiddelen aaneen tot de drie stappen uit
        \cref{prop:g8:sexual-reproduction:core}.
  \item Sluit de open dag af met één zin over waarom
        gezondheidscentra deze anatomie überhaupt onderwijzen --- kennis
        tegen verhalen, zoals \cref{ex:g8:puberty:questions} betoogde.
\end{enumerate}
\end{problem}
